% Bagian: intuitionistic-logic
% Bab: introduction
% Subbagian: natural-deduction

\documentclass[../../../include/open-logic-section]{subfiles}

\begin{document}

\olfileid[id]{int}{int}{ntd}

\olsection{Deduksi Alami}

Deduksi alami tanpa aturan $\FalseCl$ merupakan sistem !!{derivation} baku
bagi logika intuisionistik. Kita mengulangi aturan-aturannya di sini dan
menjelaskan motivasinya dengan menggunakan interpretasi BHK. Dalam setiap
kasus, kita dapat memandang suatu aturan sebagai aturan yang memungkinkan kita
menyimpulkan bahwa jika premis-premisnya mempunyai konstruksi, kesimpulannya
juga mempunyai konstruksi.

Karena !!{derivation}s deduksi alami mempunyai asumsi yang belum dilepaskan, kita
sebaiknya memandang !!a{derivation} semacam itu, misalnya derivasi $!A$ dari
asumsi !!{undischarged}~$\Gamma$, sebagai fungsi yang mengubah konstruksi bagi
semua $!B \in \Gamma$ menjadi konstruksi bagi~$!A$. Jika terdapat
!!a{derivation} bagi~$!A$ tanpa asumsi !!{undischarged}, maka terdapat
konstruksi bagi~$!A$ dalam pengertian interpretasi BHK. Namun, untuk keperluan
pembahasan, kita akan menghilangkan~$\Gamma$ jika tidak diperlukan.

Suatu asumsi~$!A$ sendiri merupakan !!a{derivation} bagi~$!A$ dari asumsi
!!{undischarged}~$!A$. Hal ini sesuai dengan interpretasi BHK: fungsi identitas
pada konstruksi mengubah setiap konstruksi bagi~$!A$ menjadi konstruksi
bagi~$!A$.

\subsection{Konjungsi}

\begin{defish}
\AxiomC{$!A$}
\AxiomC{$!B$}
\RightLabel{\Intro{\land}}
\BinaryInfC{$!A \land !B$}
\DisplayProof
\hfill
\begin{tabular}{l}
\AxiomC{$!A \land !B$}
\RightLabel{\Elim{\land}}
\UnaryInfC{$!A$}
\DisplayProof
\\[3ex]
\AxiomC{$!A \land !B$}
\RightLabel{\Elim{\land}}
\UnaryInfC{$!B$}
\DisplayProof
\end{tabular}
\end{defish}

Andaikan kita mempunyai konstruksi $N_1$, $N_2$ masing-masing bagi $!A_1$
dan $!A_2$. Dengan demikian, kita juga mempunyai konstruksi bagi $!A_1 \land
!A_2$, yakni pasangan $\tuple{N_1, N_2}$.

Menurut interpretasi BHK, konstruksi bagi $!A_1 \land !A_2$ ialah pasangan
$\tuple{N_1, N_2}$. Jadi, andaikan kita mempunyai pasangan semacam itu. Kita
pun mempunyai konstruksi bagi setiap konjung: $N_1$ merupakan konstruksi
bagi~$!A_1$, dan $N_2$ merupakan konstruksi bagi $!A_2$.

\subsection{Kondisional}

\begin{defish}
  \AxiomC{$\Discharge{!A}{u}$}
  \DeduceC{$!B$}
  \DischargeRule{$\Intro{\lif}$}{u}
  \UnaryInfC{$!A \lif !B$}
  \DisplayProof
\hfill
  \AxiomC{$!A \lif !B$}
  \AxiomC{$!A$}
  \RightLabel{$\Elim{\lif}$}
  \BinaryInfC{$!B$}
  \DisplayProof
\end{defish}

Jika kita mempunyai !!a{derivation} bagi~$!B$ dari asumsi
!!{undischarged}~$!A$, terdapat fungsi~$f$ yang mengubah konstruksi bagi~$!A$
menjadi konstruksi bagi~$!B$. Fungsi yang sama merupakan konstruksi bagi $!A
\lif !B$. Jadi, jika premis $\Intro\lif$ mempunyai konstruksi yang bergantung
pada suatu konstruksi bagi~$!A$, kesimpulan $!A \lif !B$ mempunyai konstruksi.

Sebaliknya, andaikan terdapat konstruksi~$N$ bagi $!A$ dan $f$ bagi $!A \lif
!B$. Konstruksi bagi $!A \lif !B$ ialah fungsi yang mengubah konstruksi bagi
$!A$ menjadi konstruksi bagi~$!B$. Jadi, $f(N)$ merupakan konstruksi bagi~$!B$;
yakni, kesimpulan $\Elim\lif$ mempunyai konstruksi.

\subsection{Disjungsi}

\begin{defish}
\begin{tabular}{l}
\AxiomC{$!A$}
\RightLabel{\Intro{\lor}}
\UnaryInfC{$!A \lor !B$}
\DisplayProof
\\[3ex]
\AxiomC{$!B$}
\RightLabel{\Intro{\lor}}
\UnaryInfC{$!A \lor !B$}
\DisplayProof
\end{tabular}
\hfill
\AxiomC{$!A \lor !B$}
\AxiomC{$\Discharge{!A}{n}$}
\DeduceC{$!C$}
\AxiomC{$\Discharge{!B}{n}$}
\DeduceC{$!C$}
\DischargeRule{\Elim{\lor}}{n}
\TrinaryInfC{$!C$}
\DisplayProof
\end{defish}

Jika kita mempunyai konstruksi~$N_i$ bagi~$!A_i$, kita dapat mengubahnya
menjadi konstruksi~$\tuple{i, N_i}$ bagi~$!A_1 \lor !A_2$. Sebaliknya,
andaikan kita mempunyai konstruksi bagi~$!A_1 \lor !A_2$, yakni pasangan
$\tuple{i, N_i}$ dengan $N_i$ suatu konstruksi bagi~$!A_i$, serta fungsi
$f_1$, $f_2$ yang masing-masing mengubah konstruksi bagi $!A_1$, $!A_2$
menjadi konstruksi bagi~$!C$. Maka $f_i(N_i)$ merupakan konstruksi bagi~$!C$,
yakni kesimpulan dari~$\Elim\lor$.

\subsection{Absurditas}

\begin{defish}
  \AxiomC{$\lfalse$}
  \RightLabel{\FalseInt}
  \UnaryInfC{$!A$}
  \DisplayProof
\end{defish}

Jika kita mempunyai !!a{derivation} bagi~$\lfalse$ dari asumsi
!!{undischarged}~$!B_1$, \dots, $!B_n$, maka terdapat fungsi~$f(M_1, \dots,
M_n)$ yang mengubah konstruksi bagi $!B_1$, \dots, $!B_n$ menjadi konstruksi
bagi~$\lfalse$. Karena $\lfalse$ tidak mempunyai konstruksi, konstruksi bagi
semua $!B_1$, \dots, $!B_n$ juga tidak mungkin ada. Oleh karena itu, $f$ juga
mempunyai sifat bahwa \emph{jika} $M_1$, \dots, $M_n$ masing-masing merupakan
konstruksi bagi $!B_1$, \dots, $!B_n$, \emph{maka} $f(M_1, \dots, M_n)$
merupakan konstruksi bagi~$!A$.

\subsection{Aturan untuk $\lnot$}

Karena $\lnot !A$ didefinisikan sebagai $!A \lif \lfalse$, secara ketat kita
tidak memerlukan aturan bagi~$\lnot$. Namun, seandainya kita menggunakannya,
aturan tersebut berbentuk sebagai berikut:

\begin{defish}
\AxiomC{$\Discharge{!A}{n}$}
\noLine
\DeduceC{$\lfalse$}
\DischargeRule{\Intro{\lnot}}{n}
\UnaryInfC{$\lnot !A$}
\DisplayProof
\hfill
\AxiomC{$\lnot !A$}
\AxiomC{$!A$}
\RightLabel{\Elim{\lnot}}
\BinaryInfC{$\lfalse$}
\DisplayProof
\end{defish}


\subsection{Contoh \usetoken{P}{derivation}}

\begin{enumerate}
\item $\Proves !A \lif (\lnot !A \lif \lfalse)$, yakni $\Proves !A
  \lif ((!A \lif \lfalse) \lif \lfalse)$
  \begin{prooftree}
    \AxiomC{$\Discharge{!A}{2}$}
    \AxiomC{$\Discharge{!A \lif \lfalse}{1}$}
    \RightLabel{\Elim\lif}
    \BinaryInfC{$\lfalse$}
    \DischargeRule{\Intro\lif}{1}
    \UnaryInfC{$(!A \lif \lfalse)\lif \lfalse$}
    \DischargeRule{\Intro\lif}{2}
    \UnaryInfC{$!A \lif ((!A \lif \lfalse) \lif \lfalse)$}
  \end{prooftree}

\item $\Proves ((!A \land !B) \lif !C) \lif (!A \lif (!B \lif !C))$
  \begin{prooftree}
    \AxiomC{$\Discharge{(!A \land !B) \lif !C}{3}$}
    \AxiomC{$\Discharge{!A}{2}$}
    \AxiomC{$\Discharge{!B}{1}$}
    \RightLabel{\Intro\land}
    \BinaryInfC{$!A \land !B$}
    \RightLabel{\Elim\lif}
    \BinaryInfC{$!C$}
    \DischargeRule{\Intro\lif}{1}
    \UnaryInfC{$!B \lif !C$}
    \DischargeRule{\Intro\lif}{2}
    \UnaryInfC{$!A \lif (!B \lif !C)$}
    \DischargeRule{\Intro\lif}{3}
    \UnaryInfC{$((!A \land !B) \lif !C) \lif (!A \lif (!B \lif !C))$}
  \end{prooftree}

\item $\Proves \lnot(!A \land \lnot !A)$, yakni $\Proves (!A \land (!A
  \lif \lfalse))\lif \lfalse$
  \begin{prooftree}
    \AxiomC{$\Discharge{!A \land (!A \lif \lfalse)}{1}$}
    \RightLabel{\Elim\land}
    \UnaryInfC{$!A \lif \lfalse$}
    \AxiomC{$\Discharge{!A \land (!A \lif \lfalse)}{1}$}
    \RightLabel{\Elim\land}
    \UnaryInfC{$!A$}
    \RightLabel{\Elim\lif}
    \BinaryInfC{$\lfalse$}
    \DischargeRule{\Intro\lif}{1}
    \UnaryInfC{$(!A \land (!A \lif \lfalse))\lif \lfalse$}
  \end{prooftree}

\item $\Proves \lnot\lnot(!A \lor \lnot !A)$, yakni $\Proves ((!A \lor
  (!A \lif \lfalse))\lif \lfalse)\lif \lfalse$
  \begin{prooftree}\hskip-3em
    \AxiomC{$\Discharge{(!A \lor (!A \lif \lfalse))\lif \lfalse}{2}$}
    \AxiomC{$\Discharge{(!A \lor (!A \lif \lfalse))\lif \lfalse}{2}$}
    \AxiomC{$\Discharge{!A}{1}$}
    \RightLabel{\Intro\lor}
    \UnaryInfC{$!A \lor (!A \lif \lfalse)$}
    \RightLabel{\Elim\lif}
    \BinaryInfC{$\lfalse$}
    \DischargeRule{\Intro\lif}{1}
    \UnaryInfC{$!A \lif \lfalse$}
    \RightLabel{\Intro\lor}
    \UnaryInfC{$!A \lor (!A \lif \lfalse)$}
    \RightLabel{\Elim\lif}
    \insertBetweenHyps{\hskip -4em}
    \BinaryInfC{$\lfalse$}
    \DischargeRule{\Intro\lif}{2}
    \UnaryInfC{$((!A \lor (!A \lif \lfalse))\lif \lfalse)\lif \lfalse$}
  \end{prooftree}
\end{enumerate}

\begin{prop}
  Jika $\Gamma \Proves !A$ dalam deduksi alami bagi logika intuisionistik,
  maka $\Gamma \Proves !A$ dalam logika klasik. Khususnya, jika $!A$
  merupakan teorema intuisionistik, formula itu juga merupakan teorema klasik.
\end{prop}

\begin{proof}
  Setiap aturan deduksi alami juga merupakan aturan dalam deduksi alami
  klasik, sehingga setiap !!{derivation} dalam logika intuisionistik juga
  merupakan !!a{derivation} dalam logika klasik.
\end{proof}

\begin{prob}
  Berikan !!{derivation}s dalam logika intuisionistik bagi !!{formula}s
  berikut:
  \begin{enumerate}
    \item $(\lnot !A \lor !B) \lif (!A \lif !B)$
    \item $\lnot\lnot\lnot !A \lif \lnot !A$
    \item $\lnot\lnot (!A \land !B) \liff (\lnot\lnot !A\land \lnot \lnot !B)$
    \item $\lnot (!A \lor !B) \liff (\lnot !A \land \lnot !B)$
    \item $(\lnot !A \lor \lnot !B) \lif \lnot (!A \land !B)$
    \item $\lnot\lnot ( !A \land !B) \lif  (\lnot\lnot !A \lor
    \lnot\lnot !B)$
  \end{enumerate}
\end{prob}

\end{document}
